% Standalone editable source. Compile with latexmk -pdf heide_math_v2.tex.
% No external figures, bibliography database, or shell escape required.
\documentclass[11pt,a4paper]{article}
\usepackage[T1]{fontenc}
\usepackage[utf8]{inputenc}
\usepackage{lmodern}
\usepackage[margin=25mm,headheight=14pt]{geometry}
\usepackage{amsmath,amssymb,bm}
\usepackage{booktabs,array,tabularx}
\usepackage{enumitem,microtype,xcolor,fancyhdr,needspace}
\usepackage[colorlinks=true,linkcolor=blue!45!black,citecolor=blue!45!black,
 urlcolor=blue!45!black,breaklinks=true]{hyperref}
\usepackage{bookmark}
\hypersetup{pdftitle={From Heide's Force Theorem to the VASP Trace},
 pdfauthor={DMI extraction project},pdfsubject={Practical DFT and VASP foundation for paired DMI extraction}}
\numberwithin{equation}{section}
\setcounter{tocdepth}{1}
\setlength{\parindent}{0pt}
\setlength{\parskip}{5pt plus 1pt minus 1pt}
\setlength{\emergencystretch}{2em}
\setlist{nosep,leftmargin=1.5em,topsep=5pt}
\renewcommand{\arraystretch}{1.17}
\newcolumntype{Y}{>{\raggedright\arraybackslash}X}
\pagestyle{fancy}
\fancyhf{}
\fancyhead[L]{\small DMI from DFT: the VASP trace}
\fancyhead[R]{\small Revision 2.4}
\fancyfoot[C]{\thepage}
\renewcommand{\headrulewidth}{0.3pt}
\newcommand{\vct}[1]{\bm{#1}}
\newcommand{\nh}{\hat{\vct n}}
\newcommand{\qh}{\vct q}
\newcommand{\kh}{\vct k}
\newcommand{\Nm}{N_{\mathrm m}}
\newcommand{\ii}{\mathrm i}
\newcommand{\Tr}{\operatorname{Tr}}
\newcommand{\Var}{\operatorname{Var}}
\newcommand{\code}[1]{\texttt{\detokenize{#1}}}

\begin{document}
\hypersetup{pageanchor=false}
\begin{titlepage}
\vspace*{15mm}
{\large DMI extraction pipeline\par}
\vspace{12mm}
{\LARGE\bfseries From Heide's Force Theorem\\[3pt] to the VASP Trace\par}
\vspace{8mm}
{\large Computing DMI exchange vectors\\
from paired, SOC-free spin spirals\par}
\vspace{10mm}
{\large Revision 2.4 \quad | \quad 26 September 2026\par}
\vspace{13mm}
\begin{abstract}
\noindent
The calculation has three parts. First, VASP prepares a SOC-free spin spiral
with the required sublattice directions. Second, a frozen-state PAW trace
gives its first-order spin--orbit energy. Third, differences between paired
spirals determine the bond Dzyaloshinskii--Moriya vectors.

This document derives the quantities needed for those steps and maps them
to VASP inputs and the Patch B dumps. Initial moments, constraint targets
and realized site directions are expressed in one lab-frame convention.
OSZICAR sphere moments supply the magnetic readback; VASP handles the
internal spiral representation. The bond model uses the measured unit vectors directly, without
extracting site angles or introducing an intermediate unwinding step.

The formulation covers the flat, phase-decorated and conical families used
by the implementation. It retains the checks needed to establish that the
DFT states, trace and bond model agree, without developing a general theory
of constrained magnetism or statistical inference.
\end{abstract}
\vfill
\textbf{Source and scope.} This revision keeps the lab-frame
interface and removes the redundant angular parametrizations. Only two
input controls remain: an F1 offset and an F2 tilt. Readback, constraints,
pairing and fitting use vectors. The bcc Fe tests and supplied VASP source support the interface;
they do not certify the pending BFO campaign or external DMI validation.
The pipeline readers and writers must adopt the revised contract.

\textbf{Build.} Run \code{latexmk -pdf heide_math_v2.tex}, or run
\code{pdflatex} twice. All text, equations and references are in this file.
\end{titlepage}
\hypersetup{pageanchor=true}
\pagenumbering{roman}
\tableofcontents
\clearpage
\pagenumbering{arabic}

\section{What the calculation measures}\label{sec:overview}

We seek bond vectors $\vct d_{ij}$ in a unit-spin Hamiltonian whose DMI term is
$\vct d_{ij}\cdot(\hat{\vct S}_i\times\hat{\vct S}_j)$.
The electronic calculation follows the first-order spin-spiral method of
Heide, Bihlmayer and Bl\"ugel~\cite{heide}. SOC-free Kohn--Sham states are
converged first; SOC is then inserted as a perturbation. We do not obtain
each data point by subtracting two fully self-consistent SOC total energies.

For one accepted spin state, the quantity evaluated by the pipeline is
\begin{equation}\label{eq:trace}
 \boxed{T(\nh)=
 \sum_{\mu n\kh s}w_{\kh}f^{(0)}_{n\kh}
 \,[\vct c^\mu_{n\kh s}]^\dagger C^\mu_{ss}(\nh)
 \,\vct c^\mu_{n\kh s}.}
\end{equation}
Here $\mu$ labels an atom, $n$ a band, $\kh$ a $k$ point, and $s$ an
internal-$z$ spin component. The entries of $\vct c$ are PAW projector
coefficients, as in the original companion. $C^\mu$ is the on-site SOC matrix.
The direction $\nh$ selects the DMI component to be measured; it is introduced
by rotating the operator, while the production spin states remain in the
internal-$z$ convention. The superscript $(0)$ means occupations of the
SOC-free state.

Each measurement uses two separately computed states, denoted $+$ and $-$.
The negative target reverses the wavevector and the $y$ component of
each target moment, keeping $x$ and $z$ unchanged. The data point is
\begin{equation}\label{eq:y}
 \boxed{y=\frac{T_+-T_-}{\Nm}.}
\end{equation}
$\Nm$ is the number of magnetic sites in the calculation cell, two in the
stated two-Fe setup. The traces are energies per calculation cell, so $y$ is
in eV per magnetic site. There is \emph{no additional factor of one-half}.
The often-used expression $2T_+/\Nm$ is an ideal symmetry limit, not the
production measurement rule.

The minimum chain of evidence is therefore straightforward: the input spin
state must be realized, its projections and occupations must describe that
same frozen state, and the model row must use the realized directions of both
members of the pair. An error in any of these steps can produce a precise
fit to the wrong physical problem.

Throughout, $\qh_{\rm frac}$ denotes reciprocal fractional coordinates and
$\vct r_{\rm frac}$ direct fractional coordinates:
$\qh\cdot\vct r=2\pi\qh_{\rm frac}\cdot\vct r_{\rm frac}$.
Moments have VASP's magnetic-moment units; $\hat{\vct S}$ denotes a unit
direction. The bond vectors have energy units. Moment magnitudes are not
silently absorbed into $\vct d$.

\section{The DFT state and the first-order SOC energy}\label{sec:dft}

\subsection{What is frozen}

For each target spiral, VASP solves the noncollinear SOC-free
Kohn--Sham problem,
\begin{equation}\label{eq:ks}
 H^{(0)}_{\rm KS}[\rho,\vct m]\,\Psi^{(0)}_{n\kh}
 =\varepsilon^{(0)}_{n\kh}\Psi^{(0)}_{n\kh}.
\end{equation}
The effective Hamiltonian contains the kinetic, ionic, Hartree and
exchange-correlation terms, together with the selected DFT+$U$ terms when
used. A constrained calculation also includes the potential generated by
the moment penalty described in Section~\ref{sec:constraint}.
The density, magnetization, occupations and selected electronic branch all
belong to that converged state.

The first-order SOC insertion is
\begin{equation}\label{eq:firstorder}
 \Delta E^{(1)}=
 \sum_{n\kh}w_{\kh}f^{(0)}_{n\kh}
 \langle\Psi^{(0)}_{n\kh}|H_{\rm SOC}|\Psi^{(0)}_{n\kh}\rangle.
\end{equation}
This is the derivative of the stationary electronic energy with respect to
a small SOC coupling, evaluated at zero SOC with the specified occupations.
Stationarity is why the first-order result can be evaluated on the
unperturbed orbitals. It is not a claim that the finite-SOC density or
orbitals remain unchanged to all orders.

In practice, the trace uses the SOC-free input orbitals without an
intervening SOC orbital update. For a metal or a nearly degenerate manifold,
occupations and smearing must be treated consistently; arbitrary changes of
occupations are not part of Eq.~\eqref{eq:firstorder}. At finite constraint
strength the reference is stationary for the \emph{constrained} functional.
That makes the insertion well defined, but does not by itself show that the
extracted DMI is independent of the constraining field.

The quantity in Eq.~\eqref{eq:firstorder} is an SOC expectation, not VASP's
full \code{TOTEN}. In particular, neither the full DFT energy nor the
constraint penalty is added to the trace. Higher-order SOC response is tested
by the external comparisons, not included through an ad hoc correction.

\subsection{Why a spin spiral fits in the chemical cell}

Without SOC, the electronic functional is invariant under a common spin
rotation, provided the selected additional terms and constraint targets
respect that rotation. Spin orientation is then not tied to the lattice.
This permits a translation to be accompanied by a spin rotation. VASP's
\code{LNONCOLLINEAR = T}, \code{LSPIRAL = T} representation uses
\begin{equation}\label{eq:gb}
 \Psi_{n\kh}(\vct r)=
 \begin{pmatrix}
 e^{\ii(\kh-\qh/2)\cdot\vct r}\,u_{n\kh\uparrow}(\vct r)\\
 e^{\ii(\kh+\qh/2)\cdot\vct r}\,u_{n\kh\downarrow}(\vct r)
 \end{pmatrix},
 \qquad u(\vct r+\vct R)=u(\vct r).
\end{equation}
The transverse magnetization rotates by $\qh\cdot\vct R$ between cells,
while its $z$ component is periodic. Thus the electronic functions stored
in the working representation are periodic even when the physical spin
directions are not. This is the generalized Bloch construction; it avoids
building a magnetic supercell for every \code{QSPIRAL}~\cite{spiralwiki}.

SOC generally breaks this spin-rotation freedom because it ties spin to
orbital motion in the lattice. The present method consequently prepares the
spiral with \code{LSORBIT = F} and evaluates SOC perturbatively afterward.
A conical spiral is compatible with the construction: its longitudinal
component stays periodic while its transverse component winds.

\clearpage
\section{Lab-frame spin targets and OSZICAR readback}\label{sec:frames}

\subsection{Specify vectors; keep angles only as input controls}

Use \code{SAXIS = 0 0 1} throughout the production SOC-free SCFs and
projection dumps. Here \emph{lab frame} means these fixed Cartesian spin
axes, independent of the propagation direction $\qh$.
Write $\hat{\vct S}_a$ for the unit moment on basis atom $a$ in the
calculation cell and $\hat{\vct S}_a^0$ for its requested direction.
The superscript $0$ means target, not a different frame.

Only one geometric operation is needed: rotation about the fixed $z$ axis,
\begin{equation}\label{eq:Rz}
 R_z(t)=
 \begin{pmatrix}
 \cos t&-\sin t&0\\
 \sin t& \cos t&0\\
 0&0&1
 \end{pmatrix}.
\end{equation}
Here $t$ is just the argument of the rotation, positive from $+x$ toward
$+y$; it is not an additional state variable. Trigonometric arguments
are in radians unless degrees are shown.
VASP's spiral boundary condition then fixes every translated copy:
\begin{equation}\label{eq:spin}
 \boxed{\hat{\vct S}_{\vct R a}
 =R_z(2\pi\qh_{\rm frac}\cdot\vct R_{\rm frac})\hat{\vct S}_a.}
\end{equation}
\code{QSPIRAL} fixes this advance \emph{between cells}. The moments
inside the calculation cell remain separate inputs and measured outputs.

For the existing two-Fe G-type families, the entire target recipe is
\begin{equation}\label{eq:target}
 \hat{\vct S}_a^0(\qh)=
 R_z(2\pi\qh_{\rm frac}\cdot\vct\tau_{a,\rm frac})
 \hat{\vct S}_a^0(0),
\end{equation}
with the vectors in the following table. The argument $0$ specifies the
recipe before adding the position-dependent spiral rotation; it does not
require another $q=0$ calculation.

\begin{center}
\begin{tabularx}{\textwidth}{@{}lccY@{}}
\toprule Family & $\hat{\vct S}_1^0(0)$ & $\hat{\vct S}_2^0(0)$ & Purpose\\\midrule
F0 & $(1,0,0)$ & $(-1,0,0)$ & Flat reference.\\
F1 & $(1,0,0)$ & $(-\cos\psi,-\sin\psi,0)$ & Offset Fe2 in the plane.\\
F2a & $(\sin\theta,0,\cos\theta)$ & $(0,0,-1)$ & Tilt Fe1; keep Fe2 axial.\\
F2b & $(0,0,1)$ & $(-\sin\theta,0,-\cos\theta)$ & Tilt Fe2; keep Fe1 axial.\\
\bottomrule
\end{tabularx}
\end{center}
These recipes use only two adjustable angles:
\begin{itemize}
\item $\psi$ is F1's signed extra rotation of Fe2 relative to its F0
target, positive from $+x$ toward $+y$. It is not the angle between the
two moments.
\item $\theta$ is F2's tilt away from the active site's axial reference:
from $+z$ toward $+x$ for Fe1, or from $-z$ toward $-x$ for Fe2,
before the position rotation in Eq.~\eqref{eq:target}.
At $\theta=0$ that site is axial; at $90^\circ$ it is planar.
\end{itemize}
The baseline uses $\psi=30^\circ$ or $\theta=30^\circ$ as applicable.
These are input controls only. Once the target vectors have been written,
neither angle is needed by the readback or bond fit. The reference spin
signs are already in the vectors; no separate sign factors are required.

\subsection{VASP handles the internal conversion}

Write the target vectors $\hat{\vct S}_a^0(\qh)$ into \code{M_CONSTR}
and multiply them by the chosen seed magnitudes for \code{MAGMOM}.
\textbf{Do not pre-unwind either input.} In the supplied VASP source,
\code{main.F} calls \code{COUNTER_SPIRAL} on \code{ATOMOM} before the
initial-density construction. The routine applies
$R_z(-2\pi\qh_{\rm frac}\cdot\vct\tau_{a,\rm frac})$ to each initial
moment~\cite{fetest}.
VASP therefore performs the internal compensation for \code{MAGMOM} itself.
The Fe constraint test independently supports lab-frame \code{M_CONSTR}.

For example, at $q=(1/4,0,0)$ and $\tau_2=(1/2,1/2,1/2)$, a requested
physical $+x$ seed is entered as $(M,0,0)$. VASP internally rotates it
by $-45^\circ$. Applying that rotation ourselves first would compensate
twice and seed the wrong direction. Equation~\eqref{eq:target}
constructs the physical target; it is not an instruction to undo VASP's
internal rotation.

On a magnetic restart, changing \code{MAGMOM} does not rotate the inherited
state. Keep the established continuation and near-target cold-start rules.
The common \code{SAXIS} basis and the SOC operator rotation used to select
an $x$ or $y$ DMI component are separate matters
(Section~\ref{sec:rotation})~\cite{saxis}. No moment conversion is to be
applied to \code{CPROJ}, \code{COCC} or \code{CSO}.

\subsection{Read the realized directions from OSZICAR}

Use the final complete per-ion sphere-moment block of the converged parent
SCF, in POSCAR ion order. Enable \code{I_CONSTRAINED_M = 1} and provide
\code{RWIGS} for every species; use \code{LAMBDA = 0} for unconstrained
readback and the selected positive value for constrained states. At zero
$\lambda$ the penalty energy and potential vanish. VASP recommends this
sphere readback for spirals because standard \code{LORBIT} analysis does
not account for their nonperiodic magnetization~\cite{spiralhowto}.

\begin{center}
\begin{tabularx}{\textwidth}{@{}lY@{}}
\toprule Quantity & Role in this revision\\\midrule
OSZICAR \code{M_int} & Unweighted sphere moment: primary site direction
and magnitude for geometry and the bond fit. Read as lab components.\\
OSZICAR \code{MW_int} & Smoothed sphere moment: signed constraint alignment,
penalty and field checks. Read as lab components.\\
OUTCAR projected moments & Optional diagnostics only; not a fallback for
missing sphere data and not the source of fitted spin directions.\\
OUTCAR and Patch B files & Retain settings, convergence/provenance and
the separate electronic data required by the SOC trace.\\
\bottomrule
\end{tabularx}
\end{center}
\code{LORBIT = 11} ignores \code{RWIGS} for its projector analysis;
the constraint-sphere integration still needs \code{RWIGS}~\cite{rwigs}.
The three moment definitions need not agree in magnitude. Fix the sphere
radii and compare each magnitude only with a reference of the same definition.
Do not substitute the total-cell \code{mag=} vector for the per-ion data.

Normalize each resolved \code{M_int} vector:
\begin{equation}\label{eq:real}
 \boxed{\hat{\vct S}_a=
 \frac{\vct M_a^{\rm int}}{|\vct M_a^{\rm int}|}.}
\end{equation}
That is the complete conversion required for the bond model. There is no
site-angle extraction, phase subtraction or reconstruction. An axial
moment remains an ordinary vector with its sign intact.
Check that \code{MW_int} and \code{M_int} directions agree within the
declared tolerance; a significant difference requires resolving the
local-spin definition before fitting, rather than switching readouts.


The signed realization check R0 compares these directions with the lab
targets. Family geometry may allow one common rotation about $z$ because
SOC-free energies are invariant under a common rotation about $z$.
Absolute \code{MW_int} alignment and field checks use the actual
\code{M_CONSTR} vectors, without such an offset. Never align ions
independently or use absolute dot products to hide a wrong spin sign.
Require a complete converged record; a truncated block or unresolved
moment cannot supply an accepted model row.

\subsection{Evidence and limits of the convention choice}

In two-atom bcc Fe, at $q=(1/4,0,0)$, the site phases differ by $45^\circ$.
The reported final azimuths were~\cite{fetest}:
\begin{center}
\begin{tabular}{@{}lccc@{}}
\toprule Run & \code{M_int} $(1,2)$ & \code{MW_int} $(1,2)$
 & OUTCAR $(1,2)$\\\midrule
$\lambda=0$ & $(-22.67,22.67)$ & $(-22.65,22.65)$ & $(-22.68,-22.35)$\\
$\lambda=1$, both targets $+x$ & $(-0.84,0.84)$ & $(-1.39,1.39)$ & $(-0.96,-44.06)$\\
\bottomrule
\end{tabular}
\end{center}
Angles are in degrees from $+x$ toward $+y$. The readouts differ by approximately $45^\circ$ on
site 2 and zero on site 1. The constraint drives the sphere moments toward
the two lab targets, supporting the conventions above. The $q=0$ control
also reproduced the nonspiral moments. That control alone cannot identify
frames, since the position-dependent phase vanishes.

These are convention tests, not a BFO realization or DMI validation.
In particular, the $\lambda=1$ weighted directions still miss the targets
by $1.39^\circ$. Preserve the test records and supplied source excerpts
with the executable identity before transferring their conclusions to a
different build. The earlier assumption that \code{M_CONSTR} must be
unwound is superseded; no generic winding/unwinding layer is required for
the chosen input and OSZICAR interface.

\section{Holding the target with VASP constraints}\label{sec:constraint}

\code{MAGMOM} starts the electronic calculation but does not hold its
directions during self-consistency. The F1 or F2 state could relax toward
another configuration, losing the independent information needed for the
DMI fit. We therefore supply physical target directions in \code{M_CONSTR},
increase \code{LAMBDA} only as needed, and verify the resulting OSZICAR
moments. The penalty prepares the state; it is not the SOC energy we fit.

The mode-1 constraint uses the \emph{weighted} sphere moment,
$\vct{MW}_I$, printed as \code{MW_int}. It is the integral of magnetization
with the smoothing weight that vanishes at the \code{RWIGS} boundary.
It is not interchangeable with the unweighted \code{M_int} or projected
\code{OUTCAR} moment~\cite{constraintwiki}. With unit target
$\hat{\vct S}_I^0$ from \code{M_CONSTR}, the added energy and potential
coefficient are
\begin{align}
 E_p&=\lambda\sum_I\left|\vct{MW}_I-
 \hat{\vct S}_I^0(\hat{\vct S}_I^0\cdot\vct{MW}_I)\right|^2,
 \label{eq:penalty}\\
 \vct B_I^c&=2\lambda\left[\vct{MW}_I-
 \hat{\vct S}_I^0(\hat{\vct S}_I^0\cdot\vct{MW}_I)\right].
 \label{eq:field}
\end{align}
The corresponding local potential is proportional to
$\vct B_I^c\cdot\vct\sigma$ times that same smoothing weight.
$\lambda$ is \code{LAMBDA}; this vector is a constraining-field coefficient,
not itself an angular torque.

The vector equations already explain the two initialization problems.
Parallel and antiparallel moments both have zero perpendicular component,
so mode 1 does not select their sign. For an exactly perpendicular moment,
the penalty is large but its first derivative under a small rotation is
zero: the field coefficient is parallel to the moment and supplies no
angular push. Other SCF effects can move it, but this is an avoidable
starting point.

These properties explain the prescribed initialization: continue F1 through
signed decorations such as $\pm30^\circ$, $\pm60^\circ$ and $\pm90^\circ$,
and start F2 near its requested cone rather than at a flat, orthogonal
dead point. Use the smallest member of the $\lambda=1,3,10,30$ ladder that
passes the required realization checks. A larger $\lambda$ can improve
alignment while changing the electronic state, so a small penalty alone is
not evidence of an unbiased DMI vector.
The dedicated diagnostic $\lambda$ scans still run their prescribed rungs;
the first-passing rule selects the production state.

The nominal R0 direction tolerance remains $1^\circ$. For normalized
vectors, test their signed dot product against $\cos(1^\circ)$.
For nonaxial sites, also apply that test to their \emph{normalized $xy$
projections}; this retains the phase tolerance even near the axis,
where a large in-plane error can look small in three dimensions.
Use the same common rotation, if allowed for family geometry, for all
sites. An axial target instead needs its signed axial direction checked.
A nonaxial target with unresolved transverse readback fails R0.
Never replace a signed dot product by its absolute value.

Apply the geometry checks to \code{M_int} and the absolute constraint
alignment check to \code{MW_int}, retaining both results.
The existing phase RMS diagnostic remains; a $1^\circ$ bound on every
resolved phase error already bounds its RMS by $1^\circ$.
Also require magnitude drift below $5\%$, using matching sphere-moment
definitions, and recorded $E_p<10^{-3}$ eV for constrained rows.
These are the campaign thresholds~\cite{projectcp} expressed as vector
checks, not additional fitted angles. Check the signed state even when
$E_p$ is zero.

In the Fe test, $\lambda=1$, $MW_{1y}\simeq-0.044$ and
$MW_{2y}\simeq0.044$ gave $E_p=0.003924$ eV. The vector printed under
\code{lambda*MW_perp} had $y$ entries $-0.088590$ and $0.088589$:
numerically it matches $\vct B^c=2\lambda\vct{MW}_{\perp}$ despite the
label. Its squared norm summed over both sites, divided by $4\lambda$,
reproduces the penalty. Preserve that factor convention in the reader;
the three-decimal moment output alone limits the precision of this check.

Do not add or subtract $E_p$ from Eq.~\eqref{eq:trace}. It shaped the
SOC-free reference state; the inserted operator is still SOC alone.
The $\lambda$, decoration and cross-family comparisons test whether these
states support one transferable set of DMI parameters. The F1
\code{LZEROZ} policy remains a probe decision; conical F2 states require a
nonzero longitudinal component. Zero Bi/O constraint targets also require
the specified source/binary check that they are skipped correctly.

\section{From the PAW expression to the trace}\label{sec:paw}

\subsection{Why the all-electron on-site matrix appears}

VASP represents the smooth spinor orbitals $\tilde\Psi$ and restores the
near-nucleus structure through PAW partial waves. Define the projector
coefficients and all-electron on-site SOC matrix by
\begin{equation}\label{eq:pawobjects}
 c^\mu_{i,n\kh s}=
 \langle\tilde p_i^\mu|\tilde\Psi_{n\kh s}\rangle,\qquad
 C^{\mu,ss'}_{ij}=
 \langle\phi_i^\mu|H_{\rm SOC}^{ss'}|\phi_j^\mu\rangle.
\end{equation}
Indices $i,j$ label orbital/partial-wave channels on atom $\mu$.
The PAW expectation of a local operator consists of its smooth expectation
and an all-electron-minus-pseudo on-site correction:
\begin{equation}\label{eq:paw}
 \langle H_{\rm SOC}\rangle=
 \langle\tilde\Psi|H_{\rm SOC}|\tilde\Psi\rangle+
 \sum_{\mu ijss'}c_{i,s}^{\mu *}
 \left[
 \langle\phi_i^\mu|H_{\rm SOC}^{ss'}|\phi_j^\mu\rangle-
 \langle\tilde\phi_i^\mu|H_{\rm SOC}^{ss'}|\tilde\phi_j^\mu\rangle
 \right]c_{j,s'}^\mu .
\end{equation}
Within the completeness of the one-center expansion, the smooth contribution
inside a sphere cancels its pseudo subtraction. Neglecting the remaining
SOC outside the spheres leaves the all-electron on-site matrix. This is the
approximation used by VASP~\cite{lsorbit}. Smooth or nodeless pseudo waves
do \emph{not} independently make their SOC matrix elements zero.

The usual atomic form is proportional to
$\xi(r)\,\vct L\cdot\vct\sigma$, with prefactors included in $\xi$.
The radial factor comes from the reference potential; the orbital matrices
come from angular momentum in VASP's harmonic basis. Angular momentum
couples equal-$l$ channels. No additional hand-built radial integral is
needed when the validated \code{CSO} dump already supplies these matrix
elements.

\subsection{Rotate the full operator, then select spin blocks}\label{sec:rotation}

In the $z$ spin basis, its spin structure is
\begin{equation}\label{eq:blocks}
 \vct L\cdot\vct\sigma=
 \begin{pmatrix}
 L_z&L_x-\ii L_y\\
 L_x+\ii L_y&-L_z
 \end{pmatrix}.
\end{equation}
The spin-flip blocks are generally nonzero. Hermiticity requires
$C_{\uparrow\downarrow}=C_{\downarrow\uparrow}^{\dagger}$; it does not
require either flip block to be Hermitian on its own.

The campaign measures the three Cartesian components, so no Euler
angles are needed. Use these fixed spin-basis matrices:
\begin{equation}\label{eq:rotation}
 U_z=I,\qquad
 U_x=\frac{1}{\sqrt2}\begin{pmatrix}1&-1\\1&1\end{pmatrix},
 \qquad
 U_y=\frac{1}{\sqrt2}\begin{pmatrix}1&-1\\\ii&\ii\end{pmatrix},
 \qquad C(\nh)=U_{\nh}^{\dagger}C_zU_{\nh}.
\end{equation}
Here $\nh$ is simply $\hat{\vct x}$, $\hat{\vct y}$ or $\hat{\vct z}$.
Each matrix satisfies $U_{\nh}\sigma_zU_{\nh}^{\dagger}
=\nh\cdot\vct\sigma$; a common complex phase of $U$ has no effect.
The orbital indices are unaffected. After conjugating the full operator,
retain $C_{\uparrow\uparrow}(\nh)$ and $C_{\downarrow\downarrow}(\nh)$
for Eq.~\eqref{eq:trace}. For the atomic form these are proportional to
$+\xi\,\nh\cdot\vct L$ and $-\xi\,\nh\cdot\vct L$.


The order is essential. For example, the $\sigma_xL_x$ part is entirely
off-diagonal in the original $z$ basis, but contributes a diagonal
$\sigma_zL_x$ part after the rotation selecting $x$.
Discarding flip blocks before rotation would lose the $x$ response.
The implementation must decode all four blocks, rotate, and only then
select the diagonal blocks. T3 checks the stored spin ordering and rotation
signs against genuine $x$, $y$ and $z$ operator dumps; the formula for $U$
alone cannot certify the dump convention.

\subsection{Why diagonal blocks suffice at generic \(q\)}

From Eq.~\eqref{eq:gb}, a translated projection of spin $s=\pm1$ acquires
$e^{\ii(\kh-s\qh/2)\cdot\vct R}$. In $c_s^*C_{ss'}c_{s'}$ the diagonal
blocks have no remaining translation phase. The flip blocks carry
$e^{\pm\ii\qh\cdot\vct R}$. Their lattice average is
\begin{equation}\label{eq:average}
 \frac{1}{N_R}\sum_{\vct R}e^{\pm\ii\qh\cdot\vct R}=0,
 \qquad \qh\notin\{\hbox{reciprocal lattice vectors}\}.
\end{equation}
Here $N_R$ counts calculation cells. The sum covers a complete
commensurate magnetic repeat or the infinite lattice.
It yields Eq.~\eqref{eq:trace}, including the conical
families when they obey the same generalized Bloch condition. The flip
\emph{contribution to the average} vanishes, not the stored matrix blocks.

At $q=0$ this argument does not apply. N0 instead evaluates the projected
$z$ operator on an in-plane collinear SOC-free reference. A spin-only
$180^\circ$ rotation about that collinear direction leaves the reference
density unchanged and changes the sign of $\sigma_z\xi L_z$.
Its expectation must therefore be zero. This is a separate null argument
for that reference, not a generic statement about decorated $q=0$ states.

\subsection{The equivalent density-matrix contraction}

For each atom and retained spin block, define
\begin{equation}\label{eq:rho}
 \rho^{\mu s}_{ij}=
 \sum_{n\kh}w_{\kh}f^{(0)}_{n\kh}
 c^\mu_{i,n\kh s}c^{\mu *}_{j,n\kh s},\qquad
 T=\sum_{\mu s}\Tr[C^\mu_{ss}(\nh)\rho^{\mu s}].
\end{equation}
This is just Eq.~\eqref{eq:trace} with the occupied states summed first.
It provides a useful independent check of the per-state implementation.

The stored convention matters: the specified \code{COCC} has
$\code{COCC}[i,j]=\rho[j,i]$. Thus the implementation contracts
\code{sum(CSO * COCC)} over matching orbital entries, rather than taking
\code{trace(CSO @ COCC)}. The latter inserts the wrong transpose.
A wrong transpose can still produce a real result, so the imaginary-part
gate cannot replace this check.

\section{Which VASP run supplies each quantity}\label{sec:inputs}

The separation of the material operator from the state-dependent projections
is the central Patch B input rule~\cite{projectcore}. The following objects
have different origins.

\Needspace{50mm}
\begin{center}
\begin{tabularx}{\textwidth}{@{}lY@{}}
\toprule Run & Object and relevant condition\\\midrule
Material reference & Converged SOC-free in-plane reference density/potential.\\
\code{dump_C} & Full \code{CSO}, with SOC enabled and
 \code{LSPIRAL = F}; dumped in-loop after \code{SET_DD_PAW} populates it.\\
Configuration SCF & SOC-free state at the requested $q$, decorations and
 cones; includes constraints when required.\\
\code{dump_P} & \code{CPROJ}, \code{COCC} and \code{FERWE} from one
 SOC-free \code{PRE_LOOP} capture, after \code{PROALL} and before orbital updates.\\
\bottomrule
\end{tabularx}
\end{center}

\code{COCC} is rebuilt from the same input projections and weights.
\code{FERWE.dat} is the primary occupation/weight source; the stage-2
\code{EIGENVAL} is a cross-check. Projections from the input state combined
with occupations after an orbital update would not evaluate the stated
frozen-state expectation.

Rank files carry explicit global indices. Assembly must require complete,
unique coverage of the expected $k$, band, spin and channel data.
Do not infer block or cyclic ordering from the number of MPI ranks.
Apply each occupation and $k$ weight once, with the normalization recorded
by the dump. Eigenvalues, weights and provenance establish which state was
contracted; equal array shapes alone do not.

The current dump controls use \code{ALGO = Normal}, \code{NELM = 1},
\code{LWAVE = F} and \code{LCHARG = F}, with \code{LDUMPONLY} disabled.
For \code{dump_P}, the early capture is what protects the SOC-free inputs.
For \code{dump_C}, the required operator is populated later in the normal
path. These dump stages are not substitutes for a converged parent SCF.
Keep the executable/source identity, parent \code{WAVECAR}/density,
structure, PAW datasets, $k$ mesh, smearing and capture metadata with the
measurement.

A common reference \code{CSO} is convenient but is an approximation across
different self-consistent states. Although it contains no occupation
weights, its radial factors depend on the material potential. Exact spin
rotation of a fixed operator does not prove that its radial factors are
unchanged by a phase decoration, cone or constraint. The transferability
checks must therefore test the resulting DMI, not just rotation algebra.
New-build T3 requires genuine sibling axis dumps from the same reference,
not synthetic rotations labeled as independent VASP data.

\section{Paired measurements and their checks}\label{sec:pairs}

Generate the negative target directly from the positive target vectors:
\begin{equation}\label{eq:partners}
 \qh^-=-\qh^+,\qquad
 \hat{\vct S}_a^{0,-}
 =\bigl(\hat S_{a,x}^{0,+},-\hat S_{a,y}^{0,+},\hat S_{a,z}^{0,+}\bigr).
\end{equation}
This keeps the longitudinal signs and rotating/axial sublattice assignments
fixed. With the recipes in Section~\ref{sec:frames}, it is exactly F1's
$(q,\psi)\mapsto(-q,-\psi)$ or F2's $q\mapsto-q$ at unchanged tilt.
Reversing only $q$ in the F1 recipe would retain the wrong decoration;
reversing an axial spin's polarity in F2 would change the family.
Both traces use the same measurement axis and common material operator.


Ideally the SOC-free states are related by reflection of the transverse
spin phase, with the corresponding $k\mapsto-k$ mapping. In real orbital
harmonics, the projected SOC operator is imaginary, while the paired
density matrices are complex conjugates. Their real expectations then have
opposite signs: $T_-=-T_+$. This requires matched occupations and
symmetry-related electronic branches. Matching the input targets alone
does not force two SCFs onto those branches.

Production consequently measures both $T_+$ and $T_-$, forms
Eq.~\eqref{eq:y}, and checks the raw sum $T_++T_-$ with P0. A large common
background is not acceptable merely because subtraction cancels it.
Measured-direction corrections to the model do not waive this raw-trace check.

The measured pair also sets the relevant uncertainty. If the raw trace
variances are $\sigma_+^2,\sigma_-^2$ and their covariance is $V$, then
\begin{equation}\label{eq:noise}
 \Var(y)=\frac{\sigma_+^2+\sigma_-^2-2V}{\Nm^2},\qquad
 \Var(T_++T_-)=\sigma_+^2+\sigma_-^2+2V.
\end{equation}
Use paired replicas to estimate these quantities. The difference and sum
do not share an interchangeable noise threshold. Measurements along
different axes can also share noise because they reuse projections; retain
that dependence in fit and comparison uncertainties.

A wavevector at a time-reversal-invariant point is not automatically a
null for arbitrary intracell directions. Adding a reciprocal lattice vector to $q$ leaves the
cell-translation rotations unchanged when the cell moment vectors are
held fixed. Regenerating targets with Eq.~\eqref{eq:target} need not
preserve those vectors.
Keep the specified production exclusions; do not manufacture a noise floor
from presumed TRIM nulls.

\section{From the paired energies to bond DMI vectors}\label{sec:model}

The DFT calculation supplies paired SOC energies. This section constructs
the corresponding spin-model energies and fits the unknown bond vectors
$\vct d_b$. The model must represent the same extended spiral as the
DFT trace, even though VASP stores its electronic state in one cell.

\subsection{Fix the sign and count each bond once}

The native Hamiltonian uses $-J\,\hat{\vct S}_i\cdot\hat{\vct S}_j$
for exchange and $+\vct d_{ij}\cdot(\hat{\vct S}_i\times\hat{\vct S}_j)$
for DMI, with every undirected bond counted once. A canonical bond $b$
chooses its orientation from $(0,a_b)$ to $(\vct R_b,a'_b)$.
Its displacement is
$\vct r_b=\vct R_b+\vct\tau_{a'_b}-\vct\tau_{a_b}$.
The displacement determines geometry and symmetry; the integer cell
translation $\vct R_b$ determines the additional spiral advance.
Reversing the bond reverses its displacement and $\vct d_b$.

\subsection{Why average over translated copies of a bond?}

The first-order DFT trace is an energy per calculation cell of the
\emph{extended spiral}, not the DMI energy of one selected cell.
Translating a bond rotates both endpoint spins together, and therefore
rotates their cross product. Its $z$ component stays fixed; its $x,y$
components cancel when averaged over a complete magnetic repeat
(or over the infinite lattice for an incommensurate spiral).
The average is of the cross products, not of each spin separately.
This is the same cancellation used in
Eq.~\eqref{eq:average}: $q$ must not be a reciprocal lattice vector.
It does not justify this reduction for arbitrary conical states at $q=0$.

Selecting the operator axis $\nh$ rotates the physical interpretation
of the texture's axis from $z$ to $\nh$, so the surviving component
couples to $\nh\cdot\vct d_b$. The model energy per calculation cell is
\begin{equation}\label{eq:ham}
 \boxed{\overline E_{\rm DM}(\nh)=
 \sum_{b\ {\rm canonical}}
 (\nh\cdot\vct d_b)
 \left[\hat{\vct S}_{a_b}\times
 R_z(2\pi\qh_{\rm frac}\cdot\vct R_{b,\rm frac})
 \hat{\vct S}_{a'_b}\right]_z.}
\end{equation}
Here $[\ldots]_z$ means take the $z$ component. This formula already
performs the average: no supercell or numerical lattice sum is needed.
The stored moments stay in the same $z$-axis convention for all three
operator axes. Division by $\Nm$ later converts per-cell energy to
per-magnetic-site energy.

For flat spins, the cross product is already along $z$ and the average
changes nothing. A rotating--axial bond instead has a cross product
entirely in the $xy$ plane, so its contribution averages to zero.
Bonds within the rotating sublattice can remain. That is the useful
separation supplied by F2.

\textbf{Use the cell translation $\vct R_b$, not the full displacement
$\vct r_b$, in this rotation.}
The measured vectors already specify the two intracell directions.
For an intracell bond $\vct R_b=0$, so the bracket is just
$[\hat{\vct S}_{a_b}\times\hat{\vct S}_{a'_b}]_z$.
Using $\vct r_b$ would count the intracell spiral rotation twice.
The vectors also carry the reference spin signs.


For exact partners, the exchange dot products are unchanged and the
averaged DMI contribution reverses sign. The paired prediction is
$2/\Nm$ times the positive-state canonical-bond sum. A full
directed-neighbor list counts each bond twice and instead requires
$1/\Nm$; using $2/\Nm$ with that full list would double the prediction.

\subsection{Material symmetry and independent parameters}

Symmetry relates the bond vectors. The implementation stores a small set
of independent coefficients $x_j$ and fixed propagated vectors
$\vct p_{bj}$:
\begin{equation}\label{eq:basis}
 \vct d_b=\sum_j\vct p_{bj}x_j.
\end{equation}
These vectors encode the allowed directions and relationships among bonds.
They must treat $\vct d$ as an axial vector: a spatial operation $R$
contributes $\det(R)R$, with an additional minus sign if the canonical
bond orientation is reversed. Once this basis is verified for the
production structure and cutoff, no general symmetry machinery is needed
in the energy evaluation.

The supplied $R3c$ BiFeO$_3$ (BFO) design with $R_{\max}=6$ \AA\ reports nine independent
parameters. F0 alone has rank four, even when several measurement axes
are used. F1 and the two F2 families supply independent combinations.
These counts belong to that stated structure and bond basis; a changed
POSCAR or cutoff requires checking them again~\cite{projectcore}.

\subsection{Build each row from the two realized states}

For one accepted state, normalize its OSZICAR \code{M_int} moments
with Eq.~\eqref{eq:real}. Its per-cell coefficient for parameter $j$ is
\begin{equation}\label{eq:A}
 \boxed{A_j=\sum_{b\ {\rm canonical}}
 \left[\hat{\vct S}_{a_b}\times
 R_z(2\pi\qh_{\rm frac}\cdot\vct R_{b,\rm frac})
 \hat{\vct S}_{a'_b}\right]_z
 \bigl(\nh\cdot\vct p_{bj}\bigr).}
\end{equation}
This is the complete geometric kernel for F0, F1 and F2. It uses the
measured vectors and the translation connecting the bond endpoints.
The rotation constructs a physical neighbor in another cell; it does
not change the frame of a stored moment. An exactly axial endpoint
automatically gives zero, without a special azimuth convention.


For measurement $m$, evaluate this expression separately for its positive
and negative states:
\begin{equation}\label{eq:M}
 \boxed{M_{mj}=\frac{A_j^+-A_j^-}{\Nm},
 \qquad y_m=\sum_j M_{mj}x_j.}
\end{equation}
This matches the measured difference in Eq.~\eqref{eq:y}.
Only for exact geometric partners may the row be reduced to
$2A_j^+/\Nm$. Each coefficient uses its own state's measured vectors;
the negative state does not inherit the positive state's geometry.

Using measured directions corrects small geometric departures within an
accepted family. It cannot rescue collapsed moments, wrong signs,
site-scattered phases or a different electronic state. Nor does the
unit-spin model automatically account for moment-magnitude changes.
Those are reasons to reject or separately investigate a row, not to
hide the discrepancy in the fitted $\vct d$.

\textbf{A normalization check.}
Take one intracell canonical bond with
$\hat{\vct S}_1=(1,0,0)$,
$\hat{\vct S}_2=(-\sqrt3/2,-1/2,0)$ and
$\nh\cdot\vct d=1$ meV.
The cross product has $z$ component $-1/2$, giving $-0.5$ meV per cell.
The exact partner flips the $y$ components and gives $+0.5$ meV.
With two magnetic sites, $y=(-0.5-0.5)/2=-0.5$ meV per magnetic site.
An additional one-half would be wrong. If the measured partners depart
from this ideal relation, evaluate both cross products rather than
doubling the first one.

\subsection{Fit all axes and families together}

Stack all accepted rows and solve one weighted least-squares problem,
\begin{equation}\label{eq:fit}
 \min_{\vct x}\sum_m
 \frac{(y_m-\sum_jM_{mj}x_j)^2}{\sigma_{y_m}^2}.
\end{equation}
This displayed form assumes independent rows. Where the replica data show
shared errors, use their covariance in the weighting and uncertainty
calculation. There is no need to perform three independent component fits:
the same $x_j$ generate the bond vectors seen by every axis.

Use a singular-value decomposition or QR solve, and inspect rank and
conditioning before interpreting the coefficients. More $q$ points do not
fix an exact rank deficiency of the selected family. Report the residuals,
uncertainties and held-out predictions along with the vectors.

Stability across constraint strengths, phase decorations and families is
the practical test of whether one bilinear DMI model describes these DFT
states. Compare estimates with their shared-data covariance, or refit
paired replicas consistently. A small residual on the training rows alone
does not establish the bond range or rule out a state-dependent interaction.

\subsection{One measurement pair, from inputs to a fitted row}

For an F1 pair, generate lab targets at $(q,+\psi)$ and $(-q,-\psi)$,
write their physical vectors directly to VASP, and converge both SOC-free
states using the prescribed parents. Accept their geometry from OSZICAR
\code{M_int} and their constraint alignment from \code{MW_int}.

Capture each accepted state's own projections and occupations before any
orbital update. Check coverage, weights and the equivalent density-matrix
contraction. For each measurement axis, rotate the common full SOC operator
before selecting its spin-diagonal blocks; do not rotate the recorded
moments or projection arrays as an extra frame conversion.

Record both traces, imaginary residuals and the raw pair sum. Form
$y=(T_+-T_-)/\Nm$ and build the matching row
$(A_j^+-A_j^-)/\Nm$ from the two measured geometries.
Fit all accepted rows, retaining their shared-data uncertainty.

Report the resulting bond vectors with site labels, cell translations,
directed displacements, Cartesian components, energy units and the stated
Hamiltonian sign/counting convention. Include the accepted row list,
uncertainties and symmetry basis needed to reproduce the result.

\section{Checks that make the result usable}\label{sec:checks}

Each check below addresses a specific way that the DFT state or its
interpretation can fail. Their detailed run order belongs in the probe and
runner protocols, rather than in the mathematical derivation.

\begin{description}[style=nextline,leftmargin=0pt,labelsep=0pt]
\item[C0, W0 and M1: is the frozen trace assembled correctly?]
Check the full SOC matrix structure and Hermiticity, global-index coverage,
occupation/$k$ weights and common capture provenance. Compare the per-state
and density-matrix contractions; the specified M1 tolerance is
$10^{-10}$ eV. In the real-harmonic $z$ representation, the diagonal orbital
block is imaginary and antisymmetric with zero diagonal, and the two spin
diagonal blocks have opposite signs.

\item[N0, P0 and I0: does the trace have the expected null, parity and reality?]
N0 uses the in-plane collinear $q=0$ reference and projected $z$ operator.
Its stated tolerance is $\max(5\times10^{-5}\ {\rm eV},3\sigma)$ with the
appropriate measured noise. P0 checks the raw pair sum against its own
noise. I0 bounds the absolute imaginary residual below $10^{-6}$ eV.
Reality is necessary but does not certify transpose, indexing or signs.

\item[T3 and signed R0: are the axes and states the intended ones?]
Compare rotated operators with genuine axis dumps from the same reference.
Retain the source and finite-$q$ evidence for the lab input/OSZICAR
contract, then apply the signed direction, magnitude and penalty checks of
Sections~\ref{sec:frames}
and~\ref{sec:constraint}. Validate the first F2 case before bulk production.

\item[T5--T8: is one model stable across the sampled states?]
Use the constraint-strength scan, decoration scan, cross-family comparisons
and paired replicas. The relevant convergence observable is the small
paired SOC trace, not merely a converged total energy. Tightening
\code{EDIFF} is useful only insofar as it stabilizes that observable.

\item[T9: does the extracted DMI agree with an independent physical calculation?]
Use the designated external anchors, such as matched direct $q$-SOC F0
calculations, suitable supercells or an independently converged exchange
calculation. Match the structure, electronic settings, bond range, spin
normalization, sign and bond counting before comparing.
\end{description}

For example, if an external code defines a directed-pair Hamiltonian with
$-\vct D_{ij}\cdot(\hat{\vct S}_i\times\hat{\vct S}_j)$, the present canonical
plus-sign convention gives $\vct D_{ij}=-\vct d_{ij}/2$.
This conversion follows from the two stated Hamiltonians; the term
\emph{DMI vector} by itself does not determine it.

The remaining physical assumptions are compact: SOC is treated to first
order on fixed SOC-free states and occupations; its PAW contribution is
approximated by all-electron on-site terms; one material radial operator
is reused; and a finite-range bilinear unit-spin model is transferable
across the accepted constrained families. The tests above address these
assumptions where possible. They are not consequences of successful file
parsing or a small algebraic residual.

The Fe tests and supplied source establish the convention evidence
summarized in Section~\ref{sec:frames}. They do not establish BFO
realization, the production executable's complete provenance, new T3 data
or external DMI anchors. M0 remains a same-state build/debug comparison.
Implement the revised lab-target writer, OSZICAR reader and
cell-translation bond kernel together; mixing one with the older
unwound conventions would change the fitted model.

\renewcommand{\refname}{References and implementation records}
\begin{thebibliography}{10}
\raggedright
\addcontentsline{toc}{section}{References and implementation records}
\bibitem{heide}
M. Heide, G. Bihlmayer and S. Bl\"ugel,
\emph{Describing Dzyaloshinskii--Moriya spirals from first principles},
Physica B \textbf{404}, 2678--2683 (2009).
\href{https://juser.fz-juelich.de/record/5598/files/HBB_FCM08.pdf}{Author preprint}.

\bibitem{spiralwiki}
VASP Wiki, \href{https://vasp.at/wiki/Spin_spirals}{Spin spirals}.
Generalized Bloch functions and magnetization.

\bibitem{lsorbit}
VASP Wiki, \href{https://vasp.at/wiki/LSORBIT}{LSORBIT}.
All-electron one-center SOC approximation.

\bibitem{saxis}
VASP Wiki, \href{https://vasp.at/wiki/SAXIS}{SAXIS}.
Spinor basis and common-axis rotations.

\bibitem{constraintwiki}
VASP Wiki, \href{https://vasp.at/wiki/I_CONSTRAINED_M}{I\_CONSTRAINED\_M}.
Mode-1 penalty, smoothing weight and moment readbacks.

\bibitem{spiralhowto}
VASP Wiki, \href{https://vasp.at/wiki/Spin-spiral_calculations}{Spin-spiral calculations}.
Intracell moments, intercell rotation and sphere-moment readback.

\bibitem{rwigs}
VASP Wiki, \href{https://vasp.at/wiki/RWIGS}{RWIGS}.
Separate roles in sphere integration and projected analysis.

\bibitem{fetest}
Project bcc Fe convention tests and user-supplied VASP source excerpts,
reviewed 20--25 September 2026. Two-site conventional cell, sites
$(0,0,0)$ and $(1/2,1/2,1/2)$, \code{SAXIS = 0 0 1},
$q=(1/4,0,0)$, $\lambda=0,1$, with $q=0$ controls.
Source excerpts: \code{main.F}, call to \code{COUNTER_SPIRAL} at line 1056;
\code{symmetry.F}, routine at lines 4324--4345; and \code{MRHOATO}.
Line numbers identify the supplied excerpts, not a universal VASP version.
Executable hashes and full run archives were not supplied with these excerpts.

\bibitem{projectcore}
Supplied project sources: the July 2026 \code{heide_math.pdf};
\code{math_document_revision_notes.md};
\code{decorated_spiral_implementation_rev2.md};
and \code{full_test_data_requirements.md}.
These specify the paired observable, decorated-family model, Patch B and
reported material-design ranks. This revision replaces the earlier
single-trace production rule and Patch A provenance.

\bibitem{projectcp}
Supplied current implementation records:
\code{preflight_review.md}, \code{next_steps.md},
\code{constraint_probe_v2_protocol.md}, and
\code{perlmutter_test_runner_spec.md}.
These specify signed realization, parent selection and pending validation.
The lab-frame revision explicitly replaces their older expected unwound
\code{M_CONSTR} convention and generic frame-adapter/readback instructions
with Sections~\ref{sec:frames} and~\ref{sec:model}. Other campaign
requirements remain in force. The glossary and older inventories provide
historical context, not evidence that the revised code has been implemented.
\end{thebibliography}
\end{document}
